% math4cs-hw10-matching-flow.tex

% !TEX program = xelatex
%%%%%%%%%%%%%%%%%%%%
% see http://mirrors.concertpass.com/tex-archive/macros/latex/contrib/tufte-latex/sample-handout.pdf
% for how to use tufte-handout
\documentclass[a4paper, justified]{tufte-handout}

\input{hw-preamble} % feel free to modify this file if you understand LaTeX well
%%%%%%%%%%%%%%%%%%%%
\title{计算机数学 (11-matching-flow: 匹配与网络流)}
\me{魏恒峰}{hfwei@hnu.edu.cn}{}{}
\date{2026年6月15日}
%%%%%%%%%%%%%%%%%%%%
\begin{document}
\maketitle
%%%%%%%%%%%%%%%%%%%%
\noplagiarism % PLEASE DON'T DELETE THIS LINE!
%%%%%%%%%%%%%%%%%%%%
\begin{abstract}
\end{abstract}
%%%%%%%%%%%%%%%%%%%%
\beginrequired
%%%%%%%%%%%%%%%

%%%%%%%%%%%%%%%
\begin{problem}
  设 $G = (X, Y, E)$ 是一个 $k$-正则 ($k > 0$) 二部图。
  请证明:
  \begin{enumerate}[(1)]
    \item $|X| = |Y|$;
    \item $G$ 有一个 $X$-完美匹配。
  \end{enumerate}
\end{problem}

\begin{proof}
\end{proof}
%%%%%%%%%%%%%%%

%%%%%%%%%%%%%%%
\begin{problem}
  请证明: 每个二部图 $G$ 都有一个大小 $\ge e(G)/\Delta(G)$ 的匹配~\footnote{$e(G)$表示$G$的边数。}。
  (提示: 使用 K\"{o}nig-Egerv\'{a}ry 定理。)
\end{problem}

\begin{proof}
\end{proof}
%%%%%%%%%%%%%%%

%%%%%%%%%%%%%%%
\begin{problem}
  设 $Y$ 为集合,
  $\mathcal{A} = \set{A_{1}, \dots, A_{m}}$
  为包含 $m$ 个集合的集合, 其中 $A_{i} \subseteq Y$ (对 $1 \le i \le m$)。
  $\mathcal{A}$ 的相异代表系 (System of Distinct Representatives; SDR)
  是 $Y$ 中 $m$ 个不同元素 $a_{1}, \dots, a_{m}$ 构成的集合,
  其中 $a_{i} \in A_{i}$ (对 $1 \le i \le m$)。

  \noindent 请证明: $\mathcal{A}$ 有 SDR 当且仅当
  \[
    \forall S \subseteq \set{1, \dots, m}.\;
      \big\lvert \bigcup_{i \in S} A_{i} \big\rvert \ge |S|.
  \]
\end{problem}

\begin{proof}
\end{proof}
%%%%%%%%%%%%%%%

%%%%%%%%%%%%%%%
\begin{problem}
  请给出以下网络的一个最大流与一个最小割。
  要求给出 Ford-Fulkerson Method 运行过程。
  \fig{width = 0.50\textwidth}{figs/network-flow}
\end{problem}

\begin{proof}
\end{proof}
%%%%%%%%%%%%%%%

%%%%%%%%%%%%%%%
\begin{problem}
  考虑下面的定理:
  \begin{theorem}[不能告诉你名字的某个著名定理]
    设 $G = (V, E)$ 是无向连通图, $v, w \in V$ 是不同的两个顶点。
    则 $v$, $w$ 之间的边不相交的 (edge-disjoint)~\footnote{
      设 $P_{1}$, $P_{2}$ 是两条$v$、$w$间的路径。
      如果 $P_{1}$ 与 $P_{2}$ 没有公共边,
      则 $P_{1}$、$P_{2}$ 是 $v$, $w$ 之间的边不相交的路径。
    } 路径的最大条数
    等于最小$vw$-边割集~\footnote{
      设 $F \subseteq E$ 为集。
      如果 $G$ 删除 $F$ 后, $v$ 与 $w$ 不再连通,
      则称 $F$ 是 $vw$-边割集。
    }的大小。
  \end{theorem}

  \fig{width = 0.50\textwidth}{figs/edge-disjoint}

  \begin{enumerate}[(1)]
    \item 考虑图中的 $v$, $w$ 顶点。
      请给出 $v$、$w$ 间的一个最大边不相交的路径集合。
    \item 考虑图中的 $v$, $w$ 顶点。
      请给出一个最小的 $vw$-边割集。
    \item 请使用最大流-最小割定理证明上述定理~\footnote{恭喜!你刚刚证明了图论中的一个著名定理。}。
  \end{enumerate}
\end{problem}

\begin{proof}
\end{proof}
%%%%%%%%%%%%%%%

%%%%%%%%%%%%%%%%%%%%
% 如果没有需要订正的题目，可以把这部分删掉
% \begincorrection
%%%%%%%%%%%%%%%%%%%%

%%%%%%%%%%%%%%%%%%%%
% 如果没有反馈，可以把这部分删掉
\beginfb

你可以写 (也可以发邮件)
\begin{itemize}
  \item 对课程及教师的建议与意见
  \item 教材中不理解的内容
  \item 希望深入了解的内容
  \item $\cdots$
\end{itemize}
%%%%%%%%%%%%%%%%%%%%
\end{document}